% !Mode:: "TeX:UTF-8"
%% !TEX root = ../../main.tex
%% ============================================================
%% TOPIC 1. THE CONCEPT OF A MATRIX
%% ============================================================
\mashgulot{1}{The concept of a matrix}
\begin{tarif}
A rectangular array of numbers (or, more generally, of elements of a field)
arranged in $m$ rows and $n$ columns is called a \textbf{matrix}. It is written as
\[\kf{A = \begin{pmatrix}
a_{11} & a_{12} & \cdots & a_{1n} \\
a_{21} & a_{22} & \cdots & a_{2n} \\
\vdots & \vdots & \ddots & \vdots \\
a_{m1} & a_{m2} & \cdots & a_{mn}
\end{pmatrix} = (a_{ij}), \quad i=\overline{1,m},  j=\overline{1,n}.}\]
The number $a_{ij}$, standing at the intersection of the $i$th row and the
$j$th column, is called an \textbf{entry} (element) of the matrix. The
\textbf{size} of the matrix is written $m \times n$ (read “$m$ by $n$”).
\end{tarif}

\subsection{The value of a polynomial at a matrix}
\begin{tarif}
Let $A$ be a square matrix of order $n$ and let
$f(x) = c_0 + c_1 x + c_2 x^2 + \cdots + c_m x^m$ be a polynomial.
The \textbf{nonnegative integer powers} of $A$ are defined by
\[\kf{A^0 = E, \qquad A^1 = A, \qquad A^k = \underbrace{A \cdot A \cdots A}_{k\text{ times}}.}\]
The \textbf{value of the polynomial} $f(x)$ at the matrix $A$ is then the matrix
\[\kf{f(A) = c_0 E + c_1 A + c_2 A^2 + \cdots + c_m A^m.}\]
Note that the constant term $c_0$ must be multiplied by the identity matrix:
$c_0$ is replaced by $c_0 E$.
\end{tarif}
\begin{misol}
Let $f(x) = 3x - 2x^2 + 7$ and $A = \begin{pmatrix} 1 & 2 \\ 0 & -1 \end{pmatrix}$.
We first compute $A^2$:
\[\kf{A^2 = \begin{pmatrix} 1 & 2 \\ 0 & -1 \end{pmatrix}\begin{pmatrix} 1 & 2 \\ 0 & -1 \end{pmatrix}
= \begin{pmatrix} 1 & 0 \\ 0 & 1 \end{pmatrix} = E.}\]
Hence
\[\kf{f(A) = 3A - 2A^2 + 7E =}\] \[\kf{3\begin{pmatrix} 1 & 2 \\ 0 & -1 \end{pmatrix}
- 2\begin{pmatrix} 1 & 0 \\ 0 & 1 \end{pmatrix}
+ 7\begin{pmatrix} 1 & 0 \\ 0 & 1 \end{pmatrix}
= \begin{pmatrix} 8 & 6 \\ 0 & 2 \end{pmatrix}.}\]
\end{misol}
\subsection{Types of matrices}
\begin{itemize}[leftmargin=1.3em]
\item A \textbf{square matrix} has as many rows as columns ($m=n$); the number $n$
is then called the \textbf{order} of the matrix.
\item A \textbf{row matrix} consists of a single row (its size is $1 \times n$):
$A = (a_1; a_2; \cdots; a_n)$.
\item A \textbf{column matrix (column vector)} consists of a single column (its size
is $m \times 1$): $A = \begin{pmatrix} a_1\\a_2\\ \vdots\\a_m\end{pmatrix}$.
\item The \textbf{zero matrix} $O$ is the matrix all of whose entries are zero.
\item A \textbf{diagonal matrix} is a square matrix whose entries off the main
diagonal are all zero: $a_{ij}=0$ for $i\neq j$.
\item The \textbf{identity matrix} $E$ (also denoted $I$) is the diagonal matrix
with ones on the main diagonal and zeros elsewhere.
\item An \textbf{upper (lower) triangular matrix} is a square matrix whose entries
below (above) the main diagonal are all zero.
\item A \textbf{symmetric matrix} is a square matrix satisfying $A = A^{T}$, that is,
$a_{ij} = a_{ji}$ for all $i,j$.
\item The \textbf{transpose} $A^{T}$ of a matrix $A$ is obtained by interchanging
the rows and columns of $A$: if $A=(a_{ij})_{m\times n}$, then
$A^{T}=(a_{ji})_{n\times m}$. In full,
\end{itemize}
\begin{center}
$A^{T}=\begin{pmatrix}
a_{11} & a_{21} & \cdots & a_{m1} \\
a_{12} & a_{22} & \cdots & a_{m2} \\
\vdots & \vdots & \ddots & \vdots \\
a_{1n} & a_{2n} & \cdots & a_{mn}
\end{pmatrix}$
 \end{center}
 The transpose has the following properties:
 1) $(A^{T})^{T}=A$;

2) $(\lambda A)^{T}=\lambda A^{T}$;

3) $(A+B)^{T}=A^{T}+B^{T}$;

4) $(AB)^{T}=B^{T}A^{T}$.

\vspace{3pt}
\begin{misol}
The matrix $A=\begin{pmatrix}2 & -1 & 0\\ 3 & 4 & 5\end{pmatrix}$ has size
$2\times 3$, with $a_{12}=-1$ and $a_{23}=5$. Its transpose is the $3\times 2$ matrix
\[\kf{A^{T}=\begin{pmatrix}2 & 3\\ -1 & 4\\ 0 & 5\end{pmatrix}.}\]
\end{misol}
Two matrices $A=(a_{ij})$ and $B=(b_{ij})$ are \textbf{equal} if they have the same
size and their corresponding entries coincide: $a_{ij}=b_{ij}$ for all $i,j$.
\bigskip
\subsection{Arithmetic operations on matrices}
\textbf{Addition and subtraction.}
The \textbf{sum} of two matrices $A=(a_{ij})$ and $B=(b_{ij})$ of the same size
$m\times n$ is the matrix $C=A+B=(a_{ij}+b_{ij})$ obtained by adding corresponding
entries. In the same way, the \textbf{difference} is $A-B=(a_{ij}-b_{ij})$.
Only matrices of the same size can be added or subtracted.

\vspace{3pt}
\begin{misol}
For $A=\begin{pmatrix}1 & 2\\ 3 & -1\end{pmatrix}$ and
$B=\begin{pmatrix}-2 & 0\\ 1 & 4\end{pmatrix}$ we have
\[\kf{A+B=\begin{pmatrix}-1 & 2\\ 4 & 3\end{pmatrix}, \qquad
A-B=\begin{pmatrix}3 & 2\\ 2 & -5\end{pmatrix}.}\]
\end{misol}
\textbf{Multiplication by a number.}
The product of a matrix $A=(a_{ij})$ and a number (scalar) $k$ is the matrix
$kA=(k\cdot a_{ij})$ obtained by multiplying every entry of $A$ by $k$.
\vspace{3pt}
\begin{misol}
For $A=\begin{pmatrix}1 & -2 & 3\\ 0 & 4 & -1\end{pmatrix}$ and $k=-2$,
\[\kf{kA = -2A = \begin{pmatrix}-2 & 4 & -6\\ 0 & -8 & 2\end{pmatrix}.}\]
\end{misol}
\textbf{Matrix multiplication.}
Let $A=(a_{ik})$ be an $m\times p$ matrix and $B=(b_{kj})$ a $p\times n$ matrix,
so that the number of columns of $A$ equals the number of rows of $B$.
The \textbf{product} of $A$ and $B$ is the $m\times n$ matrix
\[\kf{C = AB = (c_{ij}), \qquad c_{ij}=\sum\limits_{k=1}^{p} a_{ik}b_{kj}.}\]
In words, the $(i,j)$ entry of $C$ is obtained by multiplying the entries of the
$i$th row of $A$ by the corresponding entries of the $j$th column of $B$ and adding
the results.
\vspace{3pt}
\begin{misol}
For the $2\times2$ matrices
$A=\begin{pmatrix}1 & 2\\ 0 & -1\end{pmatrix}$ and
$B=\begin{pmatrix}3 & 1\\ 2 & 4\end{pmatrix}$,
\[\kf{AB=\begin{pmatrix}1\cdot3+2\cdot2 & 1\cdot1+2\cdot4 \\ 0\cdot3+(-1)\cdot2 & 0\cdot1+(-1)\cdot4\end{pmatrix}=}\] \[\kf{=\begin{pmatrix}7 & 9 \\ -2 & -4\end{pmatrix}.}\]
\end{misol}
\begin{misol}
Matrices of different sizes can also be multiplied, provided the inner dimensions
agree. Take the $2\times3$ matrix
$A=\begin{pmatrix}1 & 0 & 2\\ -1 & 3 & 1\end{pmatrix}$ and the $3\times2$ matrix
$B=\begin{pmatrix}3 & 1\\ 0 & -1\\ 2 & 4\end{pmatrix}$. Then
\[\kf{AB=\begin{pmatrix}1\cdot3+0\cdot0+2\cdot2 & 1\cdot1+0\cdot(-1)+2\cdot4\\
-1\cdot3+3\cdot0+1\cdot2 & -1\cdot1+3\cdot(-1)+1\cdot4\end{pmatrix}
=}\] \[\kf{=\begin{pmatrix}7 & 9\\ -1 & 0\end{pmatrix},}\]
a $2\times 2$ matrix.
\end{misol}
\begin{misol}
In general $AB \neq BA$: matrix multiplication is \emph{not commutative}.
For the matrices of the previous example, $BA$ is a $3\times3$ matrix, so it cannot
even be compared with the $2\times2$ matrix $AB$.
\end{misol}
\textbf{Transposition.}
\begin{sloppypar}
Transposing twice gives back the original matrix: $(A^{T})^{T}=A$.

\bigskip

Addition of matrices and multiplication of a matrix by a number are called the
\textbf{linear operations}. Let $A, B, C$ be matrices of the same size
$m \times n$, let $O$ be the zero matrix of that size, and let $k, l$ be arbitrary
numbers (scalars). The following properties hold.
\end{sloppypar}

\textbf{Properties of addition.}
\begin{enumerate}[label=\arabic*)]
\item $A+B=B+A$ \hfill (commutativity)
\item $(A+B)+C=A+(B+C)$ \hfill (associativity)
\item $A+O=O+A=A$ \hfill (existence of a zero matrix)
\item For every matrix $A$ there is an \textbf{additive inverse} $-A=(-1)\cdot A$
such that $A+(-A)=O$.
\end{enumerate}
\textbf{Properties of multiplication by a number.}
\begin{enumerate}[label=\arabic*)]
\item $k(A+B)=kA+kB$ \hfill (distributivity over matrix addition)
\item $(k+l)A=kA+lA$ \hfill (distributivity over addition of scalars)
\item $k(lA)=(kl)A=l(kA)$ \hfill (associativity)
\item $1\cdot A=A$, \quad $0\cdot A=O$
\item $(-1)\cdot A=-A$
\end{enumerate}
\textbf{Properties of multiplication and transposition.}
\begin{enumerate}[label=\arabic*)]
\item $(AB)C=A(BC)$ \hfill (associativity of multiplication)
\item $A(B+C)=AB+AC$, \quad $(A+B)C=AC+BC$ \hfill (distributivity)
\item $k(AB)=(kA)B=A(kB)$
\item $AE=EA=A$ \hfill ($E$ is the identity matrix of suitable order)
\item $(A+B)^{T}=A^{T}+B^{T}$
\item $(kA)^{T}=kA^{T}$
\item $(AB)^{T}=B^{T}A^{T}$ \hfill (the order of the factors is reversed)
\item $AB\neq BA$ \hfill (in general, multiplication is not commutative)
\end{enumerate}
\vspace{3pt}
\begin{misol}
Let $A=\begin{pmatrix}1 & -1\\ 2 & 0\end{pmatrix}$,
$B=\begin{pmatrix}0 & 3\\ -1 & 2\end{pmatrix}$, $k=2$ and $l=-1$. Then
\[\kf{kA+lB = 2\begin{pmatrix}1 & -1\\ 2 & 0\end{pmatrix} - \begin{pmatrix}0 & 3\\ -1 & 2\end{pmatrix}
=}\] \[\kf{= \begin{pmatrix}2 & -2\\ 4 & 0\end{pmatrix} + \begin{pmatrix}0 & -3\\ 1 & -2\end{pmatrix}
= \begin{pmatrix}2 & -5\\ 5 & -2\end{pmatrix}.}\]
With the same matrices, the reader may verify the distributive law
$k(A+B)=kA+kB$ for $k=2$.
\end{misol}
\bigskip
% =========================================================
Matrix algebra is the mathematical basis of much of modern information security:
cryptography, coding theory, network security and digital signal processing all
rely on it. We describe a few of these applications.
\subsection{Matrices in cryptography: the Hill cipher}
The \textbf{Hill cipher}, one of the classical symmetric encryption algorithms, is
built entirely on matrix operations. Each letter is assigned a number (for
instance, $A{=}0, B{=}1,\dots,Z{=}25$), and the plaintext is split into blocks of
$n$ letters; each block is written as an $n\times1$ column vector $p$. The secret
key is an invertible $n\times n$ \textbf{key matrix} $K$, and a block is encrypted
by the rule
\[\kf{c = Kp \pmod{26}.}\]
Decryption uses $p = K^{-1}c \pmod{26}$, where $K^{-1}$ is the inverse of $K$
modulo $26$. Such an inverse exists exactly when $\det K$ is coprime to $26$, that
is, $\gcd(\det K,26)=1$; a nonzero determinant alone is not enough.
\vspace{3pt}
\begin{misol}
Take the key matrix $K=\begin{pmatrix}3 & 2\\ 5 & 7\end{pmatrix}$
($\det K=11$, which is coprime to $26$) and the plaintext block
$p=\begin{pmatrix}7\\4\end{pmatrix}$ (the letters “H”, “E”). The encrypted
block is
\[\kf{c = Kp = \begin{pmatrix}3 & 2\\ 5 & 7\end{pmatrix}\begin{pmatrix}7\\4\end{pmatrix}
= \begin{pmatrix}29\\63\end{pmatrix} \equiv \begin{pmatrix}3\\11\end{pmatrix} \pmod{26},}\]
that is, “D”, “L”. The receiver computes $K^{-1} \pmod{26}$ and recovers the
plaintext as $p=K^{-1}c$. An attacker, on the other hand, has to find $K$ (or its
inverse), so the size of the key matrix and the way it is chosen directly affect
the strength of the cipher.
\end{misol}
\textbf{Coding theory: error-detecting and error-correcting codes.}
Matrices are also used to detect and correct errors that occur when data are
transmitted, for example over a network. A \textbf{generator matrix} $G$ turns a
message vector $m$ into a codeword $c = mG$. A \textbf{parity-check matrix} $H$
tests whether a received word $c$ is a codeword: this is the case when
$Hc^{T}=0$. If $Hc^{T}\neq 0$, the vector $Hc^{T}$, called the
\textbf{syndrome}, tells us where the error has occurred (this is how Hamming
codes work). Codes of this kind detect, and sometimes automatically correct,
damage that arises during transmission or storage of data, and so protect the
integrity of information, which is one of the aims of information security.

\textbf{Linear algebra in modern encryption schemes.}
\begin{sloppypar}
Part of modern public-key cryptography, notably \textbf{lattice-based
cryptography} (which includes schemes believed to resist quantum computers),
works with operations on large matrices and vectors. The secret key is a vector,
and the public key is obtained by multiplying this vector by a large random matrix
and adding some “noise”. Recovering the unknown vector from such a perturbed
matrix–vector equation is computationally very hard, and the security of the
scheme rests on this difficulty.

\textbf{Image protection and steganography.}
A digital image is a matrix of pixel values (a colour image is described by three
matrices, for the R, G and B channels).
\end{sloppypar}
\begin{itemize}[leftmargin=1.3em]
\item In \textbf{digital watermarking} and \textbf{steganography}, hidden data are
embedded in selected entries of the pixel matrix (or of a linear transform of it).
This is used to protect copyright and to transmit information covertly.
\item Methods used in image compression and encryption, such as the
\textbf{discrete cosine transform} (DCT) and the \textbf{singular value
decomposition} (SVD), are based directly on linear algebra operations on matrices.
\end{itemize}
\textbf{Network security analysis.}
The connections in a computer network can be represented as a \textbf{graph},
which in turn is described by its \textbf{adjacency matrix}. Computations with
this matrix make it possible to identify vulnerable nodes (those with many links,
which are attractive targets), to model possible attack paths such as the spread
of a virus, and, in intrusion detection systems (IDS), to represent traffic
features as a matrix and detect suspicious activity with machine learning
algorithms.

Matrix theory is thus much more than the study of tables of numbers: it provides
the mathematical foundation for encryption algorithms, error-correcting codes,
digital signatures and authentication, image protection and network analysis.
A specialist in information security therefore needs a working command of matrix
operations: addition, multiplication, inversion, computing determinants and so on.

\subsection{Computing with Python, Maple and Mathcad}
The operations of this section are available in every computer algebra system.
As an illustration we compute, for $A=\begin{pmatrix}1 & 2\\ 0 & -1\end{pmatrix}$
and $B=\begin{pmatrix}3 & 1\\ 2 & 4\end{pmatrix}$, the product $AB$, the
transpose $A^{T}$ and the value $f(A)$ of $f(x)=3x-2x^2+7$ found in the examples
above. In Python (NumPy) the operator \verb|@| denotes matrix multiplication,
whereas \verb|*| multiplies entrywise:
\begin{Verbatim}[fontsize=\footnotesize,frame=single,label={\small Python}]
import numpy as np
A = np.array([[1, 2], [0, -1]])
B = np.array([[3, 1], [2, 4]])
E = np.eye(2, dtype=int)
print(A @ B)                  # matrix product AB
print(A.T)                    # transpose
print(3*A - 2*(A @ A) + 7*E)  # f(A), f(x)=3x-2x^2+7
\end{Verbatim}
\begin{Verbatim}[fontsize=\footnotesize,frame=single,label={\small Output}]
[[ 7  9]
 [-2 -4]]
[[ 1  0]
 [ 2 -1]]
[[8 6]
 [0 2]]
\end{Verbatim}
In Maple the dot denotes the matrix product, and \verb|A^2| is the matrix power:
\begin{Verbatim}[fontsize=\footnotesize,frame=single,label={\small Maple}]
with(LinearAlgebra):
A := Matrix([[1, 2], [0, -1]]):
B := Matrix([[3, 1], [2, 4]]):
A . B;
Transpose(A);
3*A - 2*A^2 + 7*IdentityMatrix(2);
\end{Verbatim}
In Mathcad Prime the matrices are entered with the Matrix button (Ctrl+M), the
transpose is the operator \verb|A^T| (Ctrl+1), and the identity matrix is given by
\verb|identity(2)|. Writing \verb|+ 7| instead of \verb|+ 7*identity(2)| would add
$7$ to every entry, which is not the value of the polynomial.
\begin{Verbatim}[fontsize=\footnotesize,frame=single,label={\small Mathcad}]
A := | 1   2 |        B := | 3  1 |
     | 0  -1 |             | 2  4 |

A*B = |  7   9 |      A^T = | 1   0 |
      | -2  -4 |            | 2  -1 |

3*A - 2*A^2 + 7*identity(2) = | 8  6 |
                              | 0  2 |
\end{Verbatim}

\subsection{Exercises on matrix algebra}
\emph{1. Given a polynomial $f(x)$ and a square matrix $A$, find $f(A)$.}\bigskip
\begin{multicols}{2}\raggedright
\begin{enumerate}[label=\arabic*.,  leftmargin=*]
\item $f(x) = 3x - 2x^2 + 7$ and \\
$A = \begin{pmatrix} -4 & -3 \\ 4 & -3 \end{pmatrix}$.
\item $f(x) = 3x - x^2 + 1$ and \\
$A = \begin{pmatrix} -4 & -3 \\ 2 & 2 \end{pmatrix}$.
\item $f(x) = x - 2x^2 - 3$ and \\
$A = \begin{pmatrix} 4 & 2 \\ -4 & -3 \end{pmatrix}$.
\item $f(x) = 2x - x^2 + 7$ and \\
$A = \begin{pmatrix} -4 & -1 \\ -4 & 4 \end{pmatrix}$.
\item $f(x) = 2x - 3x^2 + 8$ and \\
$A = \begin{pmatrix} -2 & 4 \\ -3 & 0 \end{pmatrix}$.
\item $f(x) = 5x - 2x^2 - 2$ and \\
$A = \begin{pmatrix} -1 & 1 \\ -3 & 4 \end{pmatrix}$.
\item $f(x) = x - x^2 + 1$ and \\
$A = \begin{pmatrix} 3 & 4 \\ 2 & 1 \end{pmatrix}$.
\item $f(x) = 4x - 4x^2 + 6$ and \\
$A = \begin{pmatrix} 0 & -1 \\ -2 & -1 \end{pmatrix}$.
\item $f(x) = x - 3x^2 + 10$ and \\
$A = \begin{pmatrix} 1 & 3 \\ 0 & -3 \end{pmatrix}$.
\item $f(x) = x - 4x^2$ and \\
$A = \begin{pmatrix} 1 & -2 \\ 3 & 2 \end{pmatrix}$.
\item $f(x) = x - x^2 + 5$ and \\
$A = \begin{pmatrix} 1 & 1 \\ 3 & 3 \end{pmatrix}$.
\item $f(x) = x - x^2 + 3$ and \\
$A = \begin{pmatrix} 3 & -3 \\ -4 & 0 \end{pmatrix}$.
\item $f(x) = 5x - 4x^2 + 4$ and \\
$A = \begin{pmatrix} 2 & 1 \\ -4 & 3 \end{pmatrix}$.
\item $f(x) = 3x - 2x^2 - 2$ and \\
$A = \begin{pmatrix} 3 & -4 \\ -1 & 0 \end{pmatrix}$.
\item $f(x) = 2x - 2x^2 + 7$ and \\
$A = \begin{pmatrix} 2 & 3 \\ -3 & -2 \end{pmatrix}$.
\item $f(x) = 4x - 4x^2 + 3$ and \\
$A = \begin{pmatrix} -2 & 2 \\ 4 & 0 \end{pmatrix}$.
\item $f(x) = 4x - 3x^2 + 7$ and \\
$A = \begin{pmatrix} -1 & -2 \\ -3 & -2 \end{pmatrix}$.
\item $f(x) = 2x - 2x^2 + 2$ and \\
$A = \begin{pmatrix} -4 & 3 \\ -2 & 0 \end{pmatrix}$.
\item $f(x) = 3x - x^2 - 1$ and \\
$A = \begin{pmatrix} 2 & 4 \\ 1 & 1 \end{pmatrix}$.
\item $f(x) = 2x - x^2 + 9$ and \\
$A = \begin{pmatrix} 4 & 2 \\ 2 & 2 \end{pmatrix}$.
\item $f(x) = 4x - x^2 + 10$ and \\
$A = \begin{pmatrix} 2 & -4 \\ -1 & -3 \end{pmatrix}$.
\item $f(x) = 2x - 4x^2$ and \\
$A = \begin{pmatrix} -3 & 1 \\ -4 & -3 \end{pmatrix}$.
\item $f(x) = x - 2x^2 - 2$ and \\
$A = \begin{pmatrix} 1 & -4 \\ -3 & -1 \end{pmatrix}$.
\item $f(x) = 5x - 4x^2 - 1$ and \\
$A = \begin{pmatrix} 0 & 1 \\ 1 & 3 \end{pmatrix}$.
\item $f(x) = x - x^2 + 10$ and \\
$A = \begin{pmatrix} 3 & 3 \\ 3 & 0 \end{pmatrix}$.
\item $f(x) = x - 2x^2 - 2$ and \\
$A = \begin{pmatrix} 1 & 0 \\ 3 & -2 \end{pmatrix}$.
\item $f(x) = 5x - x^2 + 1$ and \\
$A = \begin{pmatrix} 4 & 1 \\ -2 & 4 \end{pmatrix}$.
\item $f(x) = x - 3x^2 - 3$ and \\
$A = \begin{pmatrix} 0 & 4 \\ 1 & -2 \end{pmatrix}$.
\item $f(x) = 3x - 2x^2 + 5$ and \\
$A = \begin{pmatrix} -1 & -1 \\ -1 & 2 \end{pmatrix}$.
\item $f(x) = 2x - 2x^2 + 10$ and \\
$A = \begin{pmatrix} 1 & -4 \\ -4 & 0 \end{pmatrix}$.
\end{enumerate}
\end{multicols}

\emph{2. Show that $A + A^{T}$ is a symmetric matrix.}
\begin{misol}
Show that $A+A^{T}$ is symmetric for $A=\begin{pmatrix}1 & 3 \\ -2 & 5\end{pmatrix}$.

\textbf{Solution.} Here $A^{T}=\begin{pmatrix}1 & -2 \\ 3 & 5\end{pmatrix}$, so
\[\kf{A+A^{T}=\begin{pmatrix}1 & 3 \\ -2 & 5\end{pmatrix}+\begin{pmatrix}1 & -2 \\ 3 & 5\end{pmatrix}=\begin{pmatrix}2 & 1 \\ 1 & 10\end{pmatrix}=B.}\]
Since $B^{T}=\begin{pmatrix}2 & 1 \\ 1 & 10\end{pmatrix}=B$, the matrix $A+A^{T}$ is symmetric.

In general, for any square matrix $A$,
$(A+A^{T})^{T}=A^{T}+(A^{T})^{T}=A^{T}+A=A+A^{T}$.
\end{misol}
\begin{tasks}(2)
\task $A = \begin{pmatrix} -2 & -3 & 5 \\ 3 & -1 & 1 \\ 6 & 5 & -1 \end{pmatrix}$ .
\task $A = \begin{pmatrix} -5 & -3 \\ -5 & -3 \end{pmatrix}$ .
\task $A = \begin{pmatrix} -3 & -1 & -3 \\ 1 & 3 & 3 \\ -6 & 1 & 4 \end{pmatrix}$ .
\task $A = \begin{pmatrix} 6 & 4 \\ -5 & 4 \end{pmatrix}$ .
\task $A = \begin{pmatrix} 0 & 6 \\ 5 & 6 \end{pmatrix}$ .
\task $A = \begin{pmatrix} 1 & -4 \\ 0 & 6 \end{pmatrix}$ .
\task $A = \begin{pmatrix} -5 & 6 \\ 5 & 0 \end{pmatrix}$ .
\task $A = \begin{pmatrix} 0 & 5 & -5 \\ 5 & -4 & -4 \\ -4 & -6 & -4 \end{pmatrix}$ .
\task $A = \begin{pmatrix} 1 & 6 & 4\\ -4 & 3 & 3 \\ 1 & 4 & -1 \end{pmatrix}$ .
\task $A = \begin{pmatrix} 2 & 2 \\ -4 & -6 \end{pmatrix}$ .
\task $A = \begin{pmatrix} 6 & 5 \\ 4 & -5 \end{pmatrix}$ .
\task $A = \begin{pmatrix} 5 & -4 & 0 \\ -3 & -3 & -6 \\ -2 & -3 & -2 \end{pmatrix}$ .
\task $A = \begin{pmatrix} -3 & 6 & 3 \\ -1 & -2 & 2 \\ 0 & -4 & -6 \end{pmatrix}$ .
\task $A = \begin{pmatrix} 1 & 4 \\ 3 & 2 \end{pmatrix}$ .
\task $A = \begin{pmatrix} 2 & -4 & 2 \\ -4 & 2 & 2 \\ -6 & 1 & 6 \end{pmatrix}$ .
\task $A = \begin{pmatrix} 3 & -6 \\ 6 & 6 \end{pmatrix}$ .
\task $A = \begin{pmatrix} -4 & -4 \\ 1 & 3 \end{pmatrix}$ .
\task $A = \begin{pmatrix} 2 & -6 \\ -1 & 4 \end{pmatrix}$ .
\task $A = \begin{pmatrix} 2 & 2 & 1 \\ 6 & 6 & -5 \\ 2 & -6 & -3 \end{pmatrix}$ .
\task $A = \begin{pmatrix} -2 & -6 \\ 6 & -5 \end{pmatrix}$ .
\task $A = \begin{pmatrix} 1 & 2 & -6 \\ 6 & -5 & 1 \\ -1 & 3 & 2 \end{pmatrix}$ .
\task $A = \begin{pmatrix} 2 & -3 & 5 \\ -2 & 1 & 2 \\ 2 & 6 & 1 \end{pmatrix}$ .
\task $A = \begin{pmatrix} -3 & 5 & 2 \\ -2 & 2 & -3 \\ 1 & -4 & 0 \end{pmatrix}$ .
\task $A = \begin{pmatrix} 0 & 1 \\ -1 & -5 \end{pmatrix}$ .
\task $A = \begin{pmatrix} 0 & -5 \\ -3 & 4 \end{pmatrix}$ .
\task $A = \begin{pmatrix} 6 & -5 \\ 6 & -4 \end{pmatrix}$ .
\task $A = \begin{pmatrix} -4 & -2 \\ -4 & 1 \end{pmatrix}$ .
\task $A = \begin{pmatrix} 5 & -5 \\ 0 & 1 \end{pmatrix}$ .
\task $A = \begin{pmatrix} 4 & -3 \\ -4 & 5 \end{pmatrix}$ .
\task $A = \begin{pmatrix} 2 & 0 & -1 \\ 0 & -3 & -1 \\ -1 & -5 & 5 \end{pmatrix}$ .
\end{tasks}
\emph{3. Express the vector $Z$ as a linear combination of the vectors $X$ and $Y$.}
\begin{misol}
Let $X=\begin{pmatrix}1\\2\\-1\end{pmatrix}$, $Y=\begin{pmatrix}3\\0\\4\end{pmatrix}$ and $Z=\begin{pmatrix}5\\4\\2\end{pmatrix}$. Write $Z$ in the form $Z=\alpha X+\beta Y$.

\textbf{Solution.} The vector equation $\alpha\begin{pmatrix}1\\2\\-1\end{pmatrix}+\beta\begin{pmatrix}3\\0\\4\end{pmatrix}=\begin{pmatrix}5\\4\\2\end{pmatrix}$ is equivalent to the system
\[\kf{\begin{cases}\alpha+3\beta=5,\\2\alpha=4,\\-\alpha+4\beta=2.\end{cases}}\]
The second equation gives $\alpha=2$, and then the first gives $2+3\beta=5$, so
$\beta=1$. These values also satisfy the third equation: $-2+4=2$.
Hence $Z=2X+Y$.
\end{misol}
\bigskip
\begin{enumerate}[label=\arabic*.,  leftmargin=*]
\item $X = \begin{pmatrix} 0 \\ -5 \\ 0 \end{pmatrix}$, $Y = \begin{pmatrix} 3 \\ 2 \\ 2 \end{pmatrix}$ and $Z = \begin{pmatrix} 6 \\ 24 \\ 4 \end{pmatrix}$.
\item $X = \begin{pmatrix} 0 \\ 3 \\ 4 \end{pmatrix}$, $Y = \begin{pmatrix} -1 \\ 3 \\ -4 \end{pmatrix}$ and $Z = \begin{pmatrix} 1 \\ -12 \\ -8 \end{pmatrix}$.
\item $X = \begin{pmatrix} -4 \\ -4 \\ -1 \end{pmatrix}$, $Y = \begin{pmatrix} -1 \\ -5 \\ -3 \end{pmatrix}$ and $Z = \begin{pmatrix} 6 \\ -2 \\ -4 \end{pmatrix}$.
\item $X = \begin{pmatrix} 5 \\ -1 \\ 1 \end{pmatrix}$, $Y = \begin{pmatrix} -3 \\ 3 \\ 3 \end{pmatrix}$ and $Z = \begin{pmatrix} 12 \\ 0 \\ 6 \end{pmatrix}$.
\item $X = \begin{pmatrix} -4 \\ -1 \\ -5 \end{pmatrix}$, $Y = \begin{pmatrix} -3 \\ 1 \\ -4 \end{pmatrix}$ and $Z = \begin{pmatrix} 25 \\ 1 \\ 32 \end{pmatrix}$.
\item $X = \begin{pmatrix} -1 \\ -4 \\ 4 \end{pmatrix}$, $Y = \begin{pmatrix} -2 \\ -4 \\ -1 \end{pmatrix}$ and $Z = \begin{pmatrix} -3 \\ 0 \\ -15 \end{pmatrix}$.
\item $X = \begin{pmatrix} -5 \\ 0 \\ 3 \end{pmatrix}$, $Y = \begin{pmatrix} 1 \\ -1 \\ 4 \end{pmatrix}$ and $Z = \begin{pmatrix} 6 \\ 4 \\ -22 \end{pmatrix}$.
\item $X = \begin{pmatrix} 3 \\ -2 \\ -4 \end{pmatrix}$, $Y = \begin{pmatrix} -3 \\ -1 \\ -5 \end{pmatrix}$ and $Z = \begin{pmatrix} -3 \\ 5 \\ 13 \end{pmatrix}$.
\item $X = \begin{pmatrix} -1 \\ 5 \\ -1 \end{pmatrix}$, $Y = \begin{pmatrix} 3 \\ -2 \\ -1 \end{pmatrix}$ and $Z = \begin{pmatrix} 9 \\ 7 \\ -7 \end{pmatrix}$.
\item $X = \begin{pmatrix} 5 \\ -3 \\ -1 \end{pmatrix}$, $Y = \begin{pmatrix} 0 \\ -5 \\ -1 \end{pmatrix}$ and $Z = \begin{pmatrix} -20 \\ 32 \\ 8 \end{pmatrix}$.
\item $X = \begin{pmatrix} -5 \\ 3 \\ 3 \end{pmatrix}$, $Y = \begin{pmatrix} -2 \\ 3 \\ 2 \end{pmatrix}$ and $Z = \begin{pmatrix} -1 \\ 6 \\ 3 \end{pmatrix}$.
\item $X = \begin{pmatrix} -4 \\ 5 \\ 5 \end{pmatrix}$, $Y = \begin{pmatrix} 1 \\ 5 \\ 2 \end{pmatrix}$ and $Z = \begin{pmatrix} -14 \\ 30 \\ 24 \end{pmatrix}$.
\item $X = \begin{pmatrix} 3 \\ -1 \\ -2 \end{pmatrix}$, $Y = \begin{pmatrix} -2 \\ 0 \\ -2 \end{pmatrix}$ and $Z = \begin{pmatrix} -10 \\ 2 \\ 0 \end{pmatrix}$.
\item $X = \begin{pmatrix} 0 \\ -5 \\ -3 \end{pmatrix}$, $Y = \begin{pmatrix} -5 \\ -4 \\ 5 \end{pmatrix}$ and $Z = \begin{pmatrix} 10 \\ -2 \\ -16 \end{pmatrix}$.
\item $X = \begin{pmatrix} -5 \\ -4 \\ 5 \end{pmatrix}$, $Y = \begin{pmatrix} 1 \\ 3 \\ 5 \end{pmatrix}$ and $Z = \begin{pmatrix} 1 \\ -8 \\ -25 \end{pmatrix}$.
\item $X = \begin{pmatrix} 2 \\ -3 \\ -3 \end{pmatrix}$, $Y = \begin{pmatrix} -1 \\ 2 \\ -5 \end{pmatrix}$ and $Z = \begin{pmatrix} 1 \\ -1 \\ -8 \end{pmatrix}$.
\item $X = \begin{pmatrix} 3 \\ 0 \\ -2 \end{pmatrix}$, $Y = \begin{pmatrix} -5 \\ -1 \\ -2 \end{pmatrix}$ and $Z = \begin{pmatrix} 13 \\ 2 \\ 2 \end{pmatrix}$.
\item $X = \begin{pmatrix} -5 \\ 0 \\ 1 \end{pmatrix}$, $Y = \begin{pmatrix} -4 \\ 2 \\ -1 \end{pmatrix}$ and $Z = \begin{pmatrix} -16 \\ -2 \\ 5 \end{pmatrix}$.
\item $X = \begin{pmatrix} -2 \\ 3 \\ -5 \end{pmatrix}$, $Y = \begin{pmatrix} -4 \\ -1 \\ -4 \end{pmatrix}$ and $Z = \begin{pmatrix} -4 \\ -8 \\ 2 \end{pmatrix}$.
\item $X = \begin{pmatrix} 4 \\ -5 \\ 1 \end{pmatrix}$, $Y = \begin{pmatrix} -5 \\ -1 \\ -1 \end{pmatrix}$ and $Z = \begin{pmatrix} 11 \\ 8 \\ 2 \end{pmatrix}$.
\item $X = \begin{pmatrix} 4 \\ 3 \\ -3 \end{pmatrix}$, $Y = \begin{pmatrix} 5 \\ 4 \\ 1 \end{pmatrix}$ and $Z = \begin{pmatrix} 19 \\ 15 \\ 0 \end{pmatrix}$.
\item $X = \begin{pmatrix} -3 \\ -1 \\ 4 \end{pmatrix}$, $Y = \begin{pmatrix} 5 \\ -3 \\ -5 \end{pmatrix}$ and $Z = \begin{pmatrix} -2 \\ -10 \\ 6 \end{pmatrix}$.
\item $X = \begin{pmatrix} 3 \\ -3 \\ 3 \end{pmatrix}$, $Y = \begin{pmatrix} 3 \\ 4 \\ -5 \end{pmatrix}$ and $Z = \begin{pmatrix} -12 \\ -9 \\ 12 \end{pmatrix}$.
\item $X = \begin{pmatrix} -5 \\ -5 \\ -3 \end{pmatrix}$, $Y = \begin{pmatrix} 5 \\ 0 \\ -4 \end{pmatrix}$ and $Z = \begin{pmatrix} 5 \\ -10 \\ -18 \end{pmatrix}$.
\item $X = \begin{pmatrix} 3 \\ -5 \\ 5 \end{pmatrix}$, $Y = \begin{pmatrix} -5 \\ 5 \\ 3 \end{pmatrix}$ and $Z = \begin{pmatrix} -18 \\ 20 \\ 4 \end{pmatrix}$.
\item $X = \begin{pmatrix} -1 \\ -5 \\ 2 \end{pmatrix}$, $Y = \begin{pmatrix} -4 \\ 3 \\ 3 \end{pmatrix}$ and $Z = \begin{pmatrix} -13 \\ 27 \\ 6 \end{pmatrix}$.
\item $X = \begin{pmatrix} -4 \\ 2 \\ -1 \end{pmatrix}$, $Y = \begin{pmatrix} -4 \\ -1 \\ -2 \end{pmatrix}$ and $Z = \begin{pmatrix} 8 \\ -1 \\ 3 \end{pmatrix}$.
\item $X = \begin{pmatrix} 5 \\ 2 \\ 2 \end{pmatrix}$, $Y = \begin{pmatrix} 1 \\ -4 \\ 2 \end{pmatrix}$ and $Z = \begin{pmatrix} -21 \\ -4 \\ -10 \end{pmatrix}$.
\item $X = \begin{pmatrix} -4 \\ 4 \\ -3 \end{pmatrix}$, $Y = \begin{pmatrix} 0 \\ -1 \\ 5 \end{pmatrix}$ and $Z = \begin{pmatrix} 8 \\ -4 \\ -14 \end{pmatrix}$.
\item $X = \begin{pmatrix} 2 \\ -5 \\ 2 \end{pmatrix}$, $Y = \begin{pmatrix} -1 \\ 5 \\ -4 \end{pmatrix}$ and $Z = \begin{pmatrix} -5 \\ 20 \\ -14 \end{pmatrix}$.
\end{enumerate}
\begin{misol}
Find $a$ and $b$ from $a\begin{pmatrix}2\\-1\end{pmatrix}+b\begin{pmatrix}3\\4\end{pmatrix}=\begin{pmatrix}13\\6\end{pmatrix}$.

\textbf{Solution.} Comparing the entries of $\begin{pmatrix}2a+3b\\-a+4b\end{pmatrix}=\begin{pmatrix}13\\6\end{pmatrix}$, we obtain the system

$\begin{cases}2a+3b=13,\\-a+4b=6.\end{cases}$
From the second equation $a=4b-6$. Substituting this into the first gives
$2(4b-6)+3b=13$, i.e.\ $11b=25$, so $b=\frac{25}{11}$ and
$a=4\cdot\frac{25}{11}-6=\frac{100-66}{11}=\frac{34}{11}$.
Thus $a=\frac{34}{11}$, $b=\frac{25}{11}$.
\end{misol}
\emph{4. Find the unknown numbers $a$ and $b$ from the vector equation.}
\bigskip
\begin{multicols}{2}\raggedright
\small
\begin{enumerate}[label=\arabic*.,  leftmargin=*]
\item $a\begin{pmatrix} -1 \\ 3 \end{pmatrix} + b\begin{pmatrix} -1 \\ 2 \end{pmatrix} = \begin{pmatrix} -18 \\ 45 \end{pmatrix}$.
\item $a\begin{pmatrix} -4 \\ 3 \end{pmatrix} + b\begin{pmatrix} -2 \\ -1 \end{pmatrix} = \begin{pmatrix} -30 \\ 0 \end{pmatrix}$.
\item $a\begin{pmatrix} -5 \\ -1 \end{pmatrix} + b\begin{pmatrix} 2 \\ -4 \end{pmatrix} = \begin{pmatrix} -33 \\ -33 \end{pmatrix}$.
\item $a\begin{pmatrix} 1 \\ -2 \end{pmatrix} + b\begin{pmatrix} -2 \\ -4 \end{pmatrix} = \begin{pmatrix} -5 \\ -22 \end{pmatrix}$.
\item $a\begin{pmatrix} 3 \\ -1 \end{pmatrix} + b\begin{pmatrix} 0 \\ -3 \end{pmatrix} = \begin{pmatrix} 18 \\ -15 \end{pmatrix}$.
\item $a\begin{pmatrix} 0 \\ -2 \end{pmatrix} + b\begin{pmatrix} 2 \\ 2 \end{pmatrix} = \begin{pmatrix} 4 \\ -12 \end{pmatrix}$.
\item $a\begin{pmatrix} -3 \\ -5 \end{pmatrix} + b\begin{pmatrix} 2 \\ 5 \end{pmatrix} = \begin{pmatrix} -11 \\ -5 \end{pmatrix}$.
\item $a\begin{pmatrix} -1 \\ -3 \end{pmatrix} + b\begin{pmatrix} 1 \\ 0 \end{pmatrix} = \begin{pmatrix} -1 \\ -24 \end{pmatrix}$.
\item $a\begin{pmatrix} -4 \\ 0 \end{pmatrix} + b\begin{pmatrix} -5 \\ 1 \end{pmatrix} = \begin{pmatrix} -68 \\ 8 \end{pmatrix}$.
\item $a\begin{pmatrix} -4 \\ -2 \end{pmatrix} + b\begin{pmatrix} -5 \\ -1 \end{pmatrix} = \begin{pmatrix} -59 \\ -19 \end{pmatrix}$.
\item $a\begin{pmatrix} -4 \\ 1 \end{pmatrix} + b\begin{pmatrix} 1 \\ 4 \end{pmatrix} = \begin{pmatrix} -5 \\ 31 \end{pmatrix}$.
\item $a\begin{pmatrix} 1 \\ -1 \end{pmatrix} + b\begin{pmatrix} -5 \\ -1 \end{pmatrix} = \begin{pmatrix} -7 \\ -5 \end{pmatrix}$.
\item $a\begin{pmatrix} 5 \\ -1 \end{pmatrix} + b\begin{pmatrix} 5 \\ -3 \end{pmatrix} = \begin{pmatrix} 55 \\ -23 \end{pmatrix}$.
\item $a\begin{pmatrix} 1 \\ 3 \end{pmatrix} + b\begin{pmatrix} 0 \\ -2 \end{pmatrix} = \begin{pmatrix} 7 \\ 5 \end{pmatrix}$.
\item $a\begin{pmatrix} -5 \\ 5 \end{pmatrix} + b\begin{pmatrix} 1 \\ 3 \end{pmatrix} = \begin{pmatrix} -22 \\ 34 \end{pmatrix}$.
\item $a\begin{pmatrix} -5 \\ 1 \end{pmatrix} + b\begin{pmatrix} 2 \\ 4 \end{pmatrix} = \begin{pmatrix} -8 \\ 28 \end{pmatrix}$.
\item $a\begin{pmatrix} 2 \\ -5 \end{pmatrix} + b\begin{pmatrix} 3 \\ -3 \end{pmatrix} = \begin{pmatrix} 35 \\ -47 \end{pmatrix}$.
\item $a\begin{pmatrix} 1 \\ 0 \end{pmatrix} + b\begin{pmatrix} -1 \\ -1 \end{pmatrix} = \begin{pmatrix} 0 \\ -6 \end{pmatrix}$.
\item $a\begin{pmatrix} 1 \\ 5 \end{pmatrix} + b\begin{pmatrix} -2 \\ -1 \end{pmatrix} = \begin{pmatrix} -7 \\ 37 \end{pmatrix}$.
\item $a\begin{pmatrix} -4 \\ -3 \end{pmatrix} + b\begin{pmatrix} 5 \\ -3 \end{pmatrix} = \begin{pmatrix} 13 \\ -24 \end{pmatrix}$.
\item $a\begin{pmatrix} 3 \\ 2 \end{pmatrix} + b\begin{pmatrix} 3 \\ -2 \end{pmatrix} = \begin{pmatrix} 48 \\ 4 \end{pmatrix}$.
\item $a\begin{pmatrix} 2 \\ 1 \end{pmatrix} + b\begin{pmatrix} -3 \\ 3 \end{pmatrix} = \begin{pmatrix} -5 \\ 20 \end{pmatrix}$.
\item $a\begin{pmatrix} -4 \\ -3 \end{pmatrix} + b\begin{pmatrix} 0 \\ 3 \end{pmatrix} = \begin{pmatrix} -12 \\ 12 \end{pmatrix}$.
\item $a\begin{pmatrix} -2 \\ 0 \end{pmatrix} + b\begin{pmatrix} -1 \\ 4 \end{pmatrix} = \begin{pmatrix} -12 \\ 8 \end{pmatrix}$.
\item $a\begin{pmatrix} 1 \\ 1 \end{pmatrix} + b\begin{pmatrix} 1 \\ 3 \end{pmatrix} = \begin{pmatrix} 13 \\ 29 \end{pmatrix}$.
\item $a\begin{pmatrix} -1 \\ 0 \end{pmatrix} + b\begin{pmatrix} -5 \\ 2 \end{pmatrix} = \begin{pmatrix} -41 \\ 14 \end{pmatrix}$.
\item $a\begin{pmatrix} -3 \\ 5 \end{pmatrix} + b\begin{pmatrix} 3 \\ 3 \end{pmatrix} = \begin{pmatrix} -6 \\ 34 \end{pmatrix}$.
\item $a\begin{pmatrix} -1 \\ -2 \end{pmatrix} + b\begin{pmatrix} 1\\ 1 \end{pmatrix} = \begin{pmatrix} -1 \\ -10 \end{pmatrix}$.
\item $a\begin{pmatrix} -1 \\ -5 \end{pmatrix} + b\begin{pmatrix} -3 \\ -5 \end{pmatrix} = \begin{pmatrix} -35 \\ -85 \end{pmatrix}$.
\item $a\begin{pmatrix} 4 \\ 2 \end{pmatrix} + b\begin{pmatrix} -5 \\ -4 \end{pmatrix} = \begin{pmatrix} -13 \\ -20 \end{pmatrix}$.
\end{enumerate}
\end{multicols}

\begin{misol}
A matrix has \textbf{12} entries. What sizes can it have?

\textbf{Solution.} The size $m\times n$ must satisfy $m n = 12$, so we list all
ordered pairs of positive divisors of $12$ whose product is $12$:
\[\kf{1\times 12,\quad 2\times 6,\quad 3\times 4,\quad 4\times 3,\quad 6\times 2,\quad 12\times 1.}\]
There are therefore six possible sizes.
\end{misol}
\emph{5. A matrix has $a$ entries. What sizes can it have?}
\bigskip
\begin{multicols}{4}\raggedright
\small
\begin{enumerate}[label=\arabic*.,  leftmargin=*]
\item $a=72$
\item $a=36$
\item $a=30$
\item $a=42$
\item $a=64$
\item $a=28$
\item $a=88$
\item $a=44$
\item $a=48$
\item $a=75$
\item $a=54$
\item $a=40$
\item $a=66$
\item $a=63$
\item $a=45$
\item $a=12$
\item $a=16$
\item $a=60$
\item $a=18$
\item $a=96$
\item $a=84$
\item $a=70$
\item $a=56$
\item $a=80$
\item $a=24$
\item $a=90$
\item $a=20$
\item $a=81$
\item $a=32$
\item $a=50$
\end{enumerate}
\end{multicols}

\begin{misol}
 Let $A=\begin{pmatrix}3 & -1 & 2\\ 0 & 4 & -3\end{pmatrix}$,
$B=\begin{pmatrix}1 & 2 & -1\\ -3 & 0 & 5\end{pmatrix}$, $\lambda=2$ and $\mu=-1$. Find $\lambda A+\mu B$.

\textbf{Solution.}
We have
$2A=\begin{pmatrix}6 & -2 & 4\\ 0 & 8 & -6\end{pmatrix}$ and
$(-1)B=\begin{pmatrix}-1 & -2 & 1\\ 3 & 0 & -5\end{pmatrix}$, hence
\[\kf{\lambda A+\mu B = 2A-B = \begin{pmatrix}6-1 & -2-2 & 4+1\\ 0+3 & 8+0 & -6-5\end{pmatrix}
=\begin{pmatrix}5 & -4 & 5\\ 3 & 8 & -11\end{pmatrix}.}\]
\end{misol}
\emph{6. Given matrices $A$, $B$ and numbers $\lambda, \mu$, find $\lambda A + \mu B$.}
\bigskip
\begin{enumerate}[label=\arabic*.,]
\item $A = \begin{pmatrix} 3 & 4 \\ -2 & 1 \\ -1 & -2 \end{pmatrix},  B = \begin{pmatrix} 4 & -5 \\ -5 & 3 \\ -1 & 2 \end{pmatrix},  \lambda = 1,  \mu = -1$.
\item $A = \begin{pmatrix} 3 & -2 & 3 \\ -2 & -5 & 1 \\ 5 & -1 & -5 \end{pmatrix},  B = \begin{pmatrix} -5 & -2 & 2 \\ 5 & 5 & 1 \\ -4 & -1 & -2 \end{pmatrix},  \lambda = 2,  \mu = 1$.
\item $A = \begin{pmatrix} 2 & -5 & 0 \\ 1 & 0 & 5 \end{pmatrix},  B = \begin{pmatrix} 1 & -2 & -5 \\ -1 & 3 & -4 \end{pmatrix},  \lambda = -1,  \mu = 3$.
\item $A = \begin{pmatrix} -1 & -2 & -2 \\ 2 & -2 & -1 \end{pmatrix},  B = \begin{pmatrix} -1 & -4 & 4 \\ 2 & 4 & -3 \end{pmatrix},  \lambda = -1,  \mu = 3$.
\item $A = \begin{pmatrix} 5 & -5 & 4 \\ -3 & 1 & -5 \\ -2 & -5 & 4 \end{pmatrix},  B = \begin{pmatrix} -3 & 1 & -5 \\ -5 & -3 & 1 \\ 2 & 0 & -4 \end{pmatrix},  \lambda = -3,  \mu = -2$.
\item $A = \begin{pmatrix} -2 & -3 \\ 5 & 3 \\ 2 & -5 \end{pmatrix},  B = \begin{pmatrix} -1 & 5 \\ 1 & 0 \\ 0 & 2 \end{pmatrix},  \lambda = -2,  \mu = -3$.
\item $A = \begin{pmatrix} -4 & -1 \\ -4 & 0 \end{pmatrix},  B = \begin{pmatrix} 1 & -4 \\ 3 & -2 \end{pmatrix},  \lambda = 2,  \mu = 1$.
\item $A = \begin{pmatrix} 1 & -4 \\ -5 & 2 \\ -2 & 0 \end{pmatrix},  B = \begin{pmatrix} 3 & 2 \\ -2 & 0 \\ 0 & 2 \end{pmatrix},  \lambda = -4,  \mu = 2$.
\item $A = \begin{pmatrix} 5 & 1 & -5 \\ 1 & -5 & 2 \end{pmatrix},  B = \begin{pmatrix} -4 & -5 & -1 \\ -2 & -4 & 4 \end{pmatrix},  \lambda = 1,  \mu = 1$.
\item $A = \begin{pmatrix} 0 & 4 \\ -5 & -1 \\ 0 & -1 \end{pmatrix},  B = \begin{pmatrix} -1 & -5 \\ 4 & 5 \\ -4 & -5 \end{pmatrix},  \lambda = -1,  \mu = -3$.
\item $A = \begin{pmatrix} 2 & 1 & -1 \\ 1 & 2 & -3 \\ 2 & -3 & -5 \end{pmatrix},  B = \begin{pmatrix} -1 & -3 & 4 \\ -2 & 0 & 0 \\ 2 & 0 & 4 \end{pmatrix},  \lambda = -3,  \mu = 4$.
\item $A = \begin{pmatrix} 1 & -3 & -2 \\ 1 & -4 & 5 \end{pmatrix},  B = \begin{pmatrix} -5 & 2 & 3 \\ 3 & 0 & -3 \end{pmatrix},  \lambda = 2,  \mu = -3$.
\item $A = \begin{pmatrix} -1 & 4 \\ -4 & -2 \end{pmatrix},  B = \begin{pmatrix} -4 & 1 \\ 2 & 2 \end{pmatrix},  \lambda = -2,  \mu = -1$.
\item $A = \begin{pmatrix} 1 & 2 & 4 \\ 5 & -2 & 3 \end{pmatrix},  B = \begin{pmatrix} 5 & -4 & -1 \\ -1 & -1 & 4 \end{pmatrix},  \lambda = 1,  \mu = -1$.
\item $A = \begin{pmatrix} -2 & -3 & -2 \\ -2 & -3 & -1 \\ 4 & -2 & 0 \end{pmatrix},  B = \begin{pmatrix} -4 & 1 & -1 \\ -2 & 3 & 3 \\ -2 & 5 & -4 \end{pmatrix},  \lambda = 3,  \mu = -4$.
\item $A = \begin{pmatrix} -5 & 2 \\ -2 & 2 \end{pmatrix},  B = \begin{pmatrix} 0 & -5 \\ -1 & -2 \end{pmatrix},  \lambda = -3,  \mu = -4$.
\item $A = \begin{pmatrix} 4 & 4 & -2 \\ -4 & 0 & 3 \end{pmatrix},  B = \begin{pmatrix} -3 & 2 & 4 \\ -1 & 5 & -5 \end{pmatrix},  \lambda = -3,  \mu = 1$.
\item $A = \begin{pmatrix} -5 & 0 & 0 \\ -3 & -5 & -2 \end{pmatrix},  B = \begin{pmatrix} -1 & -5 & 4 \\ 5 & -2 & -5 \end{pmatrix},  \lambda = 1,  \mu = 2$.
\item $A = \begin{pmatrix} -3 & 4 \\ -1 & -4 \\ -2 & -5 \end{pmatrix},  B = \begin{pmatrix} 2 & 3 \\ 2 & -4 \\ 1 & -4 \end{pmatrix},  \lambda = 2,  \mu = 4$.
\item $A = \begin{pmatrix} 5 & 3 & -4 \\ 5 & -3 & 1 \end{pmatrix},  B = \begin{pmatrix} -1 & 1 & -1 \\ 5 & -1 & 1 \end{pmatrix},  \lambda = -4,  \mu = 1$.
\item $A = \begin{pmatrix} 1 & -5 & 0 \\ 5 & -2 & 1 \\ 1 & -2 & -5 \end{pmatrix},  B = \begin{pmatrix} 1 & -3 & 1 \\ -4 & -4 & 1 \\ 4 & 0 & 2 \end{pmatrix},  \lambda = -2,  \mu = -2$.
\item $A = \begin{pmatrix} -5 & 3 \\ -3 & 5 \end{pmatrix},  B = \begin{pmatrix} 1 & -4 \\ 4 & 4 \end{pmatrix},  \lambda = 1,  \mu = 4$.
\item $A = \begin{pmatrix} -3 & 0 & -1 \\ -3 & 3 & -3 \end{pmatrix},  B = \begin{pmatrix} -4 & -4 & 1 \\ 2 & -2 & -1 \end{pmatrix},  \lambda = -2,  \mu = -4$.
\item $A = \begin{pmatrix} 0 & -5 & 4 \\ 5 & 1 & -4 \\ 4 & -3 & 5 \end{pmatrix},  B = \begin{pmatrix} -2 & 4 & 1 \\ 4 & -2 & 2 \\ -3 & 4 & -2 \end{pmatrix},  \lambda = -4,  \mu = 2$.
\item $A = \begin{pmatrix} 1 & 0 & -4 \\ -3 & -2 & -2 \end{pmatrix},  B = \begin{pmatrix} -5 & 3 & 5 \\ -5 & 5 & 0 \end{pmatrix},  \lambda = -3,  \mu = 2$.
\item $A = \begin{pmatrix} 3 & 5 & -1 \\ 5 & 1 & -1 \\ 4 & -2 & 1 \end{pmatrix},  B = \begin{pmatrix} 1 & 5 & 0 \\ 2 & 3 & 2 \\ -3 & -5 & -5 \end{pmatrix},  \lambda = 3,  \mu = 3$.
\item $A = \begin{pmatrix} 2 & 4 & 2 \\ -3 & 2 & 1 \end{pmatrix},  B = \begin{pmatrix} -4 & -4 & -3 \\ 0 & 1 & 0 \end{pmatrix},  \lambda = -3,  \mu = 3$.
\item $A = \begin{pmatrix} -5 & 5 \\ -3 & -4 \end{pmatrix},  B = \begin{pmatrix} 0 & 3 \\ -4 & -5 \end{pmatrix},  \lambda = 4,  \mu = 2$.
\item $A = \begin{pmatrix} -5 & -4 & 4 \\ -4 & -2 & -3 \end{pmatrix},  B = \begin{pmatrix} 2 & -1 & -3 \\ 5 & -2 & -4 \end{pmatrix},  \lambda = 1,  \mu = -2$.
\item $A = \begin{pmatrix} 4 & -1 \\ 2 & -3 \\ -1 & 3 \end{pmatrix},  B = \begin{pmatrix} 2 & -2 \\ 4 & -1 \\ 4 & 3 \end{pmatrix},  \lambda = -1,  \mu = 1$.
\end{enumerate}
\bigskip
% =========================================================
\emph{7. Applied problems.}
\begin{enumerate}[label=\arabic*.,]
\item \textbf{(Hill cipher: encryption.)} Given the key matrix
$K=\begin{pmatrix}3 & 3\\ 2 & 5\end{pmatrix}$ and the plaintext block
$p=\begin{pmatrix}9\\17\end{pmatrix}$, find the encrypted block
$c = Kp \pmod{26}$.
\item \textbf{(Hill cipher: encryption.)} Compute $c=Kp \pmod{26}$ for
$K=\begin{pmatrix}5 & 8\\ 3 & 5\end{pmatrix}$ and
$p=\begin{pmatrix}4\\19\end{pmatrix}$.
\item \textbf{(Hill cipher: encryption.)} Find the encrypted block
$c=Kp \pmod{26}$ for $K=\begin{pmatrix}7 & 8\\ 11 & 11\end{pmatrix}$ and
$p=\begin{pmatrix}2\\21\end{pmatrix}$.
\item \textbf{(Hill cipher: encryption.)} Find $c=Kp \pmod{26}$ for
$K=\begin{pmatrix}9 & 4\\ 5 & 7\end{pmatrix}$ and
$p=\begin{pmatrix}13\\4\end{pmatrix}$.
\item \textbf{(Validity of a key.)} For the matrix
$K=\begin{pmatrix}4 & 9\\ 6 & 13\end{pmatrix}$ compute the determinant and
$\gcd(\det K,26)$. Decide, with justification, whether $K$ can be used as a key
for the Hill cipher.
\item \textbf{(Validity of a key.)} Carry out the same test as in Problem~5 for
the matrix $K=\begin{pmatrix}2 & 3\\ 4 & 6\end{pmatrix}$ and state your
conclusion.
\item \textbf{(Coding theory: encoding.)} A generator matrix of the binary
$(7,4)$ Hamming code is
\[\kf{G=\begin{pmatrix}
1&0&0&0&1&1&0\\
0&1&0&0&1&0&1\\
0&0&1&0&0&1&1\\
0&0&0&1&1&1&1
\end{pmatrix}.}\]
Find the codeword $c=mG \pmod 2$ for the message vector $m=(1;0;1;1)$.
\item \textbf{(Coding theory: error detection.)} In another standard form of the
binary $(7,4)$ Hamming code, in which the $j$th column of the parity-check
matrix is the binary representation of $j$ (least significant bit in the first row), the parity-check matrix is
\[\kf{H=\begin{pmatrix}
1&0&1&0&1&0&1\\
0&1&1&0&0&1&1\\
0&0&0&1&1&1&1
\end{pmatrix}.}\]
The word $c=(1;0;1;1;1;0;1)$ has been received. Compute the syndrome
$s=Hc^{T}\pmod 2$ and use it to decide whether the received word contains an
error and, if so, in which bit.
\item \textbf{(Network security: the adjacency matrix.)}

A network with five nodes
$S_1,\dots,S_5$ has the adjacency matrix
\[\kf{A=\begin{pmatrix}
0&1&1&0&0\\
1&0&0&1&0\\
1&0&0&1&1\\
0&1&1&0&1\\
0&0&1&1&0
\end{pmatrix}.}\]
Compute $A^{2}$ and use it to find the number of paths of length $2$ between each
pair of nodes. Determine also which nodes have the most connections (and may
therefore be the most vulnerable points of the network).
\item \textbf{(A matrix polynomial.)} For the key matrix
$K=\begin{pmatrix}2 & 1\\ 1 & 1\end{pmatrix}$ compute $f(K)=K^2-3K+E$, and use
the result to show that $K^{-1}=3E-K$.
\end{enumerate}
